% The self-timed modified double-tail latch as flow/comp/subckt.py builds it: the modified double
% tail's latch (cross-coupled PMOS in series with PMOS on the first stage's drains, cross-coupled
% NMOS) with its reset NMOS (rst) driven by the first stage's drains instead of the clock.
% The preamble of every isometric drawing: standalone where installed (as docs/images/*.tex),
% else article with the page cropped to the drawing. Then Latin Modern sans and the blocks.
\IfFileExists{standalone.cls}{\documentclass[tikz,border=2mm]{standalone}\newenvironment{drawing}{}{}}{%
  \documentclass{article}\usepackage{tikz}\pagestyle{empty}\newsavebox\drawingbox
  \newenvironment{drawing}{
  \begin{lrbox}{\drawingbox}}{
    \end{lrbox}%
    \pagewidth=\dimexpr\wd\drawingbox+4mm\relax\pageheight=\dimexpr\ht\drawingbox+\dp\drawingbox+4mm\relax
\hoffset=-1in\voffset=-1in\shipout\vbox{\kern2mm\hbox{\kern2mm\usebox\drawingbox}}}}
\usepackage[T1]{fontenc}
\usepackage{lmodern}
\renewcommand{\familydefault}{\sfdefault}
\usepackage{thesis-drawings}

\usepackage{thesis-circuits}
\begin{document}
\begin{drawing}
\begin{tikzpicture}[circuit]
  \draw (-1.6,4.6) node[pmos, xscale=-1](cl){};
  \draw (1.6,4.6) node[pmos](cr){};
  \draw (-1.6,3.06) node[pmos](ol){};
  \draw (1.6,3.06) node[pmos, xscale=-1](or){};
  \draw (-1.6,0.96) node[nmos, xscale=-1](nl){};
  \draw (1.6,0.96) node[nmos](nr){};
  \draw (cl.S) node[vcc]{} (cr.S) node[vcc]{};
  \draw (nl.S) node[ground]{} (nr.S) node[ground]{};
  \draw (cl.D) -- (ol.S) (cr.D) -- (or.S);
  \draw (ol.D) -- (nl.D) (or.D) -- (nr.D);
  \draw (ol.G) node[left, sig]{$P$};
  \draw (or.G) node[right, sig]{$Q$};
  \draw (cl.G) -- (nl.G) (cr.G) -- (nr.G);
  \coordinate (X) at ($(nl.D)!0.65!(ol.D)$);
  \coordinate (Y) at ($(nr.D)!0.35!(or.D)$);
  \draw (X) -- (X -| nl.G) node[jump crossing]{} -- (X -| nr.G);
  \draw (Y) -- (Y -| nr.G) node[jump crossing]{} -- (Y -| nl.G);
  \fill (X) circle (1.6pt) (Y) circle (1.6pt) (X -| nr.G) circle (1.6pt) (Y -| nl.G) circle (1.6pt);
  \draw ($(X) - (1.5,0.77)$) node[nmos](rl){};
  \draw ($(Y) + (1.5,-0.77)$) node[nmos, xscale=-1](rr){};
  \draw (rl.S) node[ground]{} (rr.S) node[ground]{};
  \draw (rl.D) -- (X) (rr.D) -- (Y);
  \draw (rl.G) node[left, sig]{$P$};
  \draw (rr.G) node[right, sig]{$Q$};
  \noteL{cl}{on}
  \noteR{ol}{on}
  \noteL{nl}{init}
  \noteR{rl}{rst}
  \draw (X) node[above left, sig]{$X$};
  \draw (Y) node[below right, sig]{$Y$};
\end{tikzpicture}
\end{drawing}
\end{document}
